% =====================================================================
% Clase -- Física 11° -- Trimestre I -- Semana 04
% Sesiones 006 (lun 21 sep., 11:00, 100 min) y 007 (mié 23 sep., 10:10, 100 min)
% Tema 2 de la guía: La carga eléctrica · hilo «De Tales al pararrayos»
% Material del docente (no se entrega a los estudiantes).
% =====================================================================
\documentclass[11pt]{article}
\makeatletter\def\input@path{{../../../../../plantillas/estilo/}}\makeatother
\usepackage{mmcantillo}
\guiadelestudiante{../../guia-didactica/guia-periodo-I-electrostatica}

\tipodocumento{Clase}
\asignatura{Física}
\titulo{Semana 4: inducción, el electroscopio y el laboratorio de signos}
\grado{11°}
\periodo{I}
% DBA: naturales grado 11 · DBA 2 — Comprende que la interacción de las cargas en reposo
%      genera fuerzas eléctricas y que cuando las cargas están en movimiento genera fuerzas
%      magnéticas.
% Evidencia trabajada: (1) identifica el tipo de carga eléctrica (positiva o negativa) que
%      adquiere un material cuando se somete a procedimientos de fricción o contacto.
% Estándar (10°–11°, «me aproximo al conocimiento como científico natural»): el laboratorio de
%      la sesión 007 (predicción, observación, registro y contraste).

\begin{document}

\encabezado*

\begin{center}
{\renewcommand{\arraystretch}{1.3}
\begin{tabularx}{\linewidth}{@{}l l l X l@{}}
  \toprule
  \textbf{Sesión} & \textbf{Fecha} & \textbf{Min} & \textbf{Subtema} & \textbf{Entregables}\\
  \midrule
  006 & lun 21 sep., 11:00 & 100 & Electrización por inducción; el electroscopio & Taller 4\\
  007 & mié 23 sep., 10:10 & 100 & Laboratorio: qué carga adquiere un material al frotarlo o al tocarlo & Tarea 4: informe (para el lun 28 sep.)\\
  \bottomrule
\end{tabularx}}
\end{center}

Referencia: \refguia{tema:carga} (apartado «Tres maneras de electrizar un cuerpo» y la serie
triboeléctrica).

% ---------------------------------------------------------------------
\seccion{Sesión 006 — lunes 21 de septiembre · 11:00 · 100 min}

\textbf{Objetivo.} Explicar la electrización por inducción, con y sin conexión a tierra, y
usar el electroscopio para detectar carga y reconocer su signo.

\begin{center}
{\renewcommand{\arraystretch}{1.3}
\begin{tabularx}{\linewidth}{@{}l l X@{}}
  \toprule
  \textbf{Momento} & \textbf{Min} & \textbf{Actividad}\\
  \midrule
  Tarea 3 & 0–15 & Revisión en parejas; en el tablero, el punto 1 (el signo depende del par)
    y el reparto del punto 2 ($\num{4e13}$ entre dos: $\num{2e13}$ cada una).\\
  Inducción & 15–40 & Retoma la pregunta del miércoles: el tubo de PVC cargado atrae
    papelitos neutros. Dibuja una esfera metálica aislada con el tubo negativo cerca: los
    electrones libres se alejan, la cara cercana queda $+$ y la lejana $-$; la carga total
    sigue en cero, pero la atracción sobre la cara cercana gana porque está más cerca. Con
    conexión a tierra: se toca la esfera con el dedo mientras el tubo está cerca, se retira
    el dedo y \emph{después} el tubo: la esfera queda con signo contrario al del inductor.\\
  Electroscopio & 40–55 & Demostración con un electroscopio: al acercar el tubo, las hojas se
    abren (inducción, sin tocar); al tocarlo, quedan abiertas (contacto, mismo signo). Si ya
    está cargado negativamente y se acerca un cuerpo negativo, las hojas se abren más; si es
    positivo, se cierran. Así se reconoce un signo, y así se usará en el laboratorio.\\
  Construcción & 55–70 & Por grupos de 4, arman el electroscopio casero (materiales abajo) y
    lo prueban con el tubo. Lo guardan para el miércoles.\\
  Taller 4 & 70–95 & En parejas.\\
  Cierre & 95–100 & ¿Por qué el tubo y la lana deben estar secos? Se pide que cada grupo
    traiga el miércoles un globo y una bolsa plástica.\\
  \bottomrule
\end{tabularx}}
\end{center}

\begin{nota}
\textbf{Electroscopio casero (uno por grupo):} un frasco de vidrio con tapa plástica, un
alambre de cobre rígido (o un clip grande estirado) de unos 12 cm que atraviesa la tapa, una
bolita de papel aluminio en el extremo de arriba y, abajo, doblado en L, dos tiras de papel
aluminio delgado (aluminio de cocina, $\qty{1}{cm} \times \qty{3}{cm}$) colgadas una junto a
la otra. Se sella el orificio de la tapa con plastilina. Lleva un frasco armado para la
demostración y el material para el resto de los grupos.
\end{nota}

% ---------------------------------------------------------------------
\seccion{Sesión 007 — miércoles 23 de septiembre · 10:10 · 100 min}

\textbf{Objetivo.} Determinar experimentalmente el signo que adquiere un material al
frotarlo con otro o al tocar un cuerpo cargado, y contrastar el resultado con la predicción
hecha con la serie triboeléctrica. (Gilbert frotó decenas de materiales y anotó cuáles
atraían; nosotros además preguntamos qué signo toman.)

\begin{center}
{\renewcommand{\arraystretch}{1.3}
\begin{tabularx}{\linewidth}{@{}l l X@{}}
  \toprule
  \textbf{Momento} & \textbf{Min} & \textbf{Actividad}\\
  \midrule
  Predicción & 0–10 & Cada grupo copia la tabla de la Tarea 4 y escribe, con la serie de
    la guía, el signo que espera para cada material.\\
  Referencia & 10–20 & Se carga el electroscopio negativamente: PVC frotado con lana y se
    toca la bolita (las hojas quedan abiertas). Comprobación: al acercar de nuevo el PVC
    frotado, las hojas se abren más.\\
  Controles & 20–25 & Se acerca cada material \emph{sin frotar}: debe verse poco o ningún
    cambio. Si las hojas se cierran un poco con un objeto neutro, anotarlo: por eso solo una
    apertura clara es una prueba concluyente.\\
  Frotamiento & 25–60 & Para cada par de la tabla: se frotan unas diez veces, se separan y se
    acerca cada material (sin tocar) a la bolita del electroscopio. Se abren más: negativo; se
    cierran: positivo. Tres repeticiones por material; se descarga el electroscopio tocándolo
    con el dedo y se vuelve a cargar si hace falta.\\
  Contacto & 60–75 & Péndulo eléctrico: se toca la bolita de aluminio con el PVC frotado y se
    acerca de nuevo el PVC (repulsión: mismo signo). Luego se acerca el vidrio frotado con la
    bolsa (atracción). Se anota el signo que tomó la bolita.\\
  Puesta en común & 75–92 & Cada grupo dicta sus resultados; se cuentan los que coinciden
    con la predicción. ¿Hubo algún par que no coincidió? Posibles causas: humedad, superficies
    sucias, poco frotamiento, carga que se va por la mano.\\
  Cierre & 92–100 & Orden y guardado del material. Se asigna la Tarea 4 (informe) para el
    lunes 28; la puesta en común completa es la sesión 008.\\
  \bottomrule
\end{tabularx}}
\end{center}

\begin{nota}
\textbf{Materiales para el laboratorio} (grupos de 4; cantidades por grupo):
\begin{itemize}
  \item el electroscopio casero de la sesión 006;
  \item un tubo de PVC de $\frac{1}{2}$ pulgada y unos 30 cm;
  \item un trozo de lana (bufanda, suéter o paño de lana natural);
  \item un vaso o frasco de vidrio liso y seco;
  \item una bolsa de polietileno (bolsa plástica de supermercado);
  \item dos hojas de papel;
  \item un globo inflado (para el par libre: globo y pelo);
  \item péndulo eléctrico: una bolita de papel aluminio de 1 cm colgada de 20 cm de hilo de
        coser, sujeto con cinta a una regla o un lápiz apoyado entre dos libros;
  \item cinta adhesiva, tijeras, plastilina y papel toalla.
\end{itemize}
Para el docente: un secador de pelo para secar los materiales (con la humedad de Cali es lo
que más afecta los resultados), un electroscopio de repuesto y lana y PVC adicionales. Se
trabaja lejos de ventiladores y de las ventanas abiertas. Seguridad: el vidrio se manipula con
cuidado; si se quiebra, no se recoge con la mano.
\end{nota}

\end{document}
